\documentclass[main]{subfiles}

\begin{document}
% ref: claim_3_2 (proof)
% title: Untitled
% migrated from Section3_1/tex_proofs/claim_3_2_AcyclicIffTopologicalOrder.tex

% claim_3_2 — standalone TeX proof.
% source: lecture-notes/lecture_notes/graphs.tex, lines 221-245
% (the LN provides a \begin{proof} immediately after the claim).
%
% This file lifts that proof essentially verbatim, with two light edits for
% self-containment:
%   (a) explicit citations of `def_3_6` (acyclic), `def_3_8` (topological
%       order), `def_3_5` (parents), `def_3_4` (directed walk);
%   (b) the LN's parenthetical "otherwise, following parent edges would
%       produce a directed cycle" is expanded into an explicit pigeonhole
%       argument, since the downstream Lean worker needs every step.

\begin{claimmark}
\begin{Lem}
        A CDMG  $G=(J,V,E,L)$ is acyclic if and only if it has a topological order.
\end{Lem}
\end{claimmark}

\noindent\textbf{Strategy.} Two directions: ($\Leftarrow$) a directed cycle
would yield a strictly decreasing chain in the linear order, which is
impossible; ($\Rightarrow$) construct the order by repeatedly extracting a
parent-less node from the not-yet-selected vertices, using acyclicity (via
pigeonhole on the finite vertex set) to guarantee such a node always exists.

\begin{proof}
Throughout, "acyclic'' means in the sense of `def_3_6': every directed walk
$p$ from a vertex $v$ to itself is the trivial (length-$0$) walk at $v$.
A \emph{topological order} of $G$, by `def_3_8', is a total order $<$ on
$J \cup V$ such that for all $v, w \in J \cup V$,
\[
  v \in \Pa^G(w) \;\Longrightarrow\; v < w,
\]
where $\Pa^G(w)$ denotes the set of parents of $w$ in $G$ (`def_3_5'), i.e.\
those $v$ with $v \tuh w \in G$ (a directed edge from $v$ to $w$).

\medskip
\noindent\textbf{($\Leftarrow$) Topological order $\Rightarrow$ acyclic.}
Suppose $<$ is a topological order of $G$, i.e.\
$v \in \Pa^G(w) \Rightarrow v < w$.

We must show every directed walk from a vertex $v$ to itself is trivial.
Suppose for contradiction that there is a non-trivial directed walk
(a directed walk in the sense of `def_3_4', i.e.\ traversing only directed
edges $\tuh$)
\[
  v = v_0 \tuh v_1 \tuh v_2 \tuh \cdots \tuh v_n = v
\]
with $n \ge 1$. Each step $v_{i-1} \tuh v_i$ means $v_{i-1} \in \Pa^G(v_i)$
(by `def_3_5'), so by the topological-order property
\[
  v_0 < v_1 < v_2 < \cdots < v_n.
\]
But $v_n = v_0$, so $v_0 < v_0$, contradicting irreflexivity of the strict
order $<$. Hence no such non-trivial walk exists, and $G$ is acyclic.

\medskip
\noindent\textbf{($\Rightarrow$) Acyclic $\Rightarrow$ topological order.}
Suppose $G$ is acyclic. We construct a topological order by repeatedly
selecting a node with no parents in the not-yet-selected subset.

Set $N := J \cup V$ and $K := |N|$ (finite since both $J$ and $V$ are finite
by `def_3_1'). We define a sequence $v_1, v_2, \dots, v_K$ inductively.
Given $v_1, \dots, v_{i-1}$, let
\[
  R_i \;:=\; N \setminus \{v_1, \dots, v_{i-1}\}
\]
be the set of remaining (not-yet-selected) nodes; note $R_1 = N$. For a
node $w \in R_i$, write
\[
  \Pa^{R_i}(w) \;:=\; \Pa^G(w) \cap R_i
\]
for the parents of $w$ that still lie in $R_i$.

\medskip\noindent\emph{Claim.} For every $i$ with $R_i \neq \emptyset$, there
exists $v \in R_i$ with $\Pa^{R_i}(v) = \emptyset$.

\smallskip\noindent\emph{Proof of claim.} Suppose for contradiction that every
$v \in R_i$ has a parent in $R_i$. Pick any $w_0 \in R_i$ and inductively
pick $w_{k+1} \in \Pa^{R_i}(w_k)$ for $k = 0, 1, 2, \dots$ (always possible
by assumption). This produces an infinite sequence $w_0, w_1, w_2, \dots$
of elements of $R_i$ with $w_{k+1} \tuh w_k \in G$ for every $k$.

Since $R_i \subseteq N$ is finite, by the pigeonhole principle there exist
indices $0 \le k < \ell$ with $w_k = w_\ell$. Then the edges
\[
  w_\ell \tuh w_{\ell-1} \tuh w_{\ell-2} \tuh \cdots \tuh w_{k+1} \tuh w_k
\]
form a directed walk from $w_\ell$ to $w_k = w_\ell$ of length $\ell - k \ge 1$,
hence non-trivial. This contradicts acyclicity of $G$ (`def_3_6'), proving the
claim. \hfill$\diamond$

\medskip
By the claim, at each step $i = 1, 2, \dots, K$ we may pick some
$v_i \in R_i$ with $\Pa^{R_i}(v_i) = \emptyset$. After $K$ steps every node
of $N$ has been chosen exactly once, so $v_1, v_2, \dots, v_K$ is an
enumeration of $J \cup V$.

Define the total order $<$ on $N$ by $v_j < v_i$ iff $j < i$
(i.e.\ the order in which nodes were extracted). This is a strict total order
because the $v_i$ are pairwise distinct.

It remains to verify the topological-order property: if $v, w \in N$ and
$v \in \Pa^G(w)$, then $v < w$. Write $w = v_i$ and $v = v_j$. Since $v$ is
a parent of $w$, in particular $v \in \Pa^G(v_i)$.

Consider the step at which $v_i$ was selected: at that moment $v_i \in R_i$
and $\Pa^{R_i}(v_i) = \emptyset$, meaning $\Pa^G(v_i) \cap R_i = \emptyset$.
Hence $v \notin R_i$, i.e.\ $v$ had already been selected, i.e.\ $v = v_j$ with
$j < i$. By definition of $<$, this means $v < w$, as required.

Thus $<$ is a topological order of $G$, and the $\Rightarrow$ direction is
proven.
\qed
\end{proof}

\end{document}
